%! AP Statistics — Unit 4, Lesson 4.3 — Homework
\documentclass[10pt]{article}
\usepackage{apstats-article}
\usepackage{apstats-boxes}

\begin{document}

\pageheader{Unit 4, Lesson 4.3}{Homework}
\namedateperiod

\vspace{0.1in}

\begin{practicebox}
\small

\textbf{1.} A bag contains 5 red, 4 blue, and 1 yellow marble. One marble is drawn
at random.

\begin{enumerate}[label=(\alph*), leftmargin=2em, itemsep=8pt]
  \item How many total outcomes are in the sample space? \blank{2cm}
  \item Calculate each probability:
        $P(\text{red}) = \blank{2cm}$ \qquad
        $P(\text{blue}) = \blank{2cm}$ \qquad
        $P(\text{yellow}) = \blank{2cm}$
  \item Verify: do the three probabilities sum to 1? Show the sum. \blank{4cm}
  \item $P(\text{not blue}) = \blank{2cm}$ using the complement rule.
  \item Interpret $P(\text{red}) = 0.5$ in a sentence using ``in the long run.''\\
        \writeline
\end{enumerate}

\vspace{8pt}
\textbf{2.} A health clinic records the blood type of each patient. In the past, the
relative frequencies have been: O = 0.44, A = 0.42, B = 0.10, AB = 0.04.

\begin{enumerate}[label=(\alph*), leftmargin=2em, itemsep=8pt]
  \item Verify these relative frequencies sum to 1. \blank{5cm}
  \item A patient is chosen at random. Calculate:
        $P(\text{not type O}) = \blank{3cm}$ \quad $P(\text{type A or type B}) = \blank{3cm}$
  \item Interpret $P(\text{type O}) = 0.44$ in the context of this clinic.\\
        \writeline\writeline
\end{enumerate}

\vspace{8pt}
\textbf{3.} A standard six-sided die is rolled twice. Use the $6 \times 6$ sample
space (36 equally likely outcomes).

\begin{enumerate}[label=(\alph*), leftmargin=2em, itemsep=8pt]
  \item $P(\text{both dice show the same number}) = \blank{3cm}$
  \item $P(\text{sum equals 2 or 12}) = \blank{3cm}$
  \item $P(\text{sum is NOT 2 or 12}) = \blank{3cm}$
  \item Is $P(\text{sum} = 7) + P(\text{sum} \ne 7) = 1$? Verify using your answers
        from the guided notes. \blank{5cm}
\end{enumerate}

\vspace{8pt}
\textbf{4.} (AP-Style Free Response) A bag holds 4 green, 3 yellow, and 1 purple
marble. One marble is drawn at random.

\begin{enumerate}[label=(\alph*), leftmargin=2em, itemsep=8pt]
  \item Calculate $P(\text{green})$ and $P(\text{not green})$. Show your work.
        \vspace{16pt}
  \item A student claims: ``Since purple has only 1 marble, $P(\text{purple}) = 0$
        because it almost never happens.'' Explain the error in this claim.\\
        \writeline\writeline
  \item Write a long-run interpretation of $P(\text{yellow})$.\\
        \writeline\writeline
\end{enumerate}

\vspace{8pt}
\textbf{5.} (AP-Style Multiple Choice) A spinner has 5 equal sections: 2 red,
2 blue, 1 green. Which is the probability of NOT spinning red?

\begin{enumerate}[label=(\Alph*), leftmargin=2em, itemsep=4pt]
  \item $2/5$
  \item $2/3$
  \item $3/5$
  \item $1/5$
\end{enumerate}

\vspace{8pt}
\textbf{Extension / Preview:} Design a spinner with 6 sections where
$P(\text{red}) = 1/3$ and $P(\text{blue}) = 1/2$. What color fills the remaining
section(s)? What is $P(\text{remaining color})$? (\textit{Preview: Lesson 4.4 explores
conditional probability --- how does knowing one event occurred change the probability
of another?})

\end{practicebox}

\end{document}
